\documentclass[10pt,a4paper]{amsart}
\usepackage[margin=19mm]{geometry}
\usepackage{amsmath,amssymb,mathtools}
\usepackage[scr=rsfs,cal=euler]{mathalfa}
\usepackage{xfrac}
\usepackage[unicode,hidelinks,bookmarks=false]{hyperref}
\newcommand{\ie}{i.e.\ }
\newcommand{\cf}{cf.\ }
\newcommand{\resp}{respectively\ }
\newcommand{\Subsection}{\subsection}
\providecommand{\nobreakdash}{\nobreak\hbox{-}}
\setlength{\emergencystretch}{2em}
\multlinegap0pt
\usepackage{amssymb}
\usepackage{enumerate}
\usepackage{xcolor}
\newtheorem{proposition}{Proposition}
\newcommand{\Acal}{\mathcal A}
\newcommand{\Bcal}{\mathcal B}
\newcommand{\Ecal}{\mathcal E}
\newcommand{\Tcal}{\mathcal T}
\newcommand{\dd}{\,\mathrm{d}}
\newcommand{\f}{\frac}
\newcommand{\p}{\partial}
\title{Global small data radial symmetric solutions of 3D semilinear Euler-Poisson-Darboux equations}
\author{Qianqian Li, Qiao Xin, Huicheng Yin}
\date{Source: arXiv:2608.29258v1 (CC BY 4.0)}
\begin{document}
\maketitle

\noindent\emph{Source and licence.} Global small data radial symmetric solutions of 3D semilinear Euler-Poisson-Darboux equations. Version arXiv:2608.29258v1 (CC BY 4.0), distributed by arXiv under the Creative Commons Attribution 4.0 International licence, http://creativecommons.org/licenses/by/4.0/, as recorded in the arXiv licence metadata for this exact version and shown on the abstract page https://arxiv.org/abs/2608.29258v1. Full licence notice: \texttt{licenses/arxiv-2608.29258-CC-BY-4.0.txt}.\par\smallskip

\noindent\textit{Editorial scope.} Counted unit: the article's proposition prop:nonlinear (source0.tex lines 2314--2349). The article's complete original proof of it is delivered with it (source0.tex lines 2351--3031, one \texttt{proof} environment of 681 lines). No mathematics is written, repaired or re-derived: the statement and the proof are the article's own lines, in the article's own notation, and no retained line is substituted or re-flowed.  Retained uncounted as the source-local context the proof needs: lines 1368--1378; lines 1843--1848; lines 1879--1896; lines 1898--1921; lines 2027--2031; lines 2037--2046; lines 2048--2068; lines 2229--2235; lines 2295--2301.  Nothing was omitted from any retained span, so the package carries no omission record.  No retained line contains a citation, so the wrapper carries no bibliography and no bibliography entry.  No retained line contains a figure environment, an image-inclusion command or drawing code, so no image is delivered and the package declares no asset dependency.  Editorial formatting only, in addition: an AMS \texttt{amsart} preamble that restates the article's own declarations and macros used by the retained lines, namely the environments \texttt{proposition} on separate counters, together with the standard packages those lines need.  The retained environments print in this document as Proposition 1 in order of appearance, whereas the article prints the same items with the numbering of its own full document.  The complete native source of the article as distributed in the arXiv e-print for this exact version is bundled unchanged as \texttt{sources/208836/source0.tex}.
\begin{equation}\label{eq:Xkappa}
\begin{aligned}
\|w\|_{X_\kappa}
=\sup_{\substack{t\geq1\\0\leq r\leq t}}
\big\{
&\Acal(t,r)\Bcal(t,r)^{\kappa-1}|w(t,r)|+
\f{\Acal(t,r)\Bcal(t,r)^{\kappa-1}}{1+r}
\left|\p_r(rw)(t,r)\right|
\big\}.
\end{aligned}
\end{equation}

\begin{equation}
\label{eq:kappa-range}
\f{2-\beta}{p-1}
<\kappa
<\min\{p+\beta-1,\,2-\lambda\}.
\end{equation}

\begin{equation}
\label{qq:40}
\begin{aligned}
c_\lambda^{-1}r(\Tcal w)(t,r)
={}&
\int_1^t
\int_{|t-s-r|}^{t-s+r}
\Ecal_\lambda(t,s,r-q)
G_w(s,q)\dd q\dd s\\
&+\int_1^{t-r}
\int_0^{t-s-r}
\left[
\Ecal_\lambda(t,s,r-q)
-\Ecal_\lambda(t,s,r+q)
\right]
G_w(s,q)\dd q\dd s.
\end{aligned}
\end{equation}

\begin{equation}
\label{qq:42}
\begin{aligned}
c_\lambda^{-1}\p_r(r\Tcal w)(t,r)
={}&
4^\lambda
\int_1^t
\left[
G_w(s,t-s+r)
+G_w(s,t-s-r)
\right]\dd s\\
&+\int_1^t
\int_{|t-s-r|}^{t-s+r}
\p_y\Ecal_\lambda(t,s,r-q)
G_w(s,q)\dd q\dd s\\
&+\int_1^{t-r}
\int_0^{t-s-r}
\left[
\p_y\Ecal_\lambda(t,s,r-q)
-\p_y\Ecal_\lambda(t,s,r+q)
\right]
G_w(s,q)\dd q\dd s.
\end{aligned}
\end{equation}

\begin{equation}
\label{qq:44}
A_0=t+r,\quad B_0=t-r,\quad
\xi=s+q,\quad \eta=q-s,
\end{equation}

\begin{equation}
\label{qq:45}
|\Ecal_\lambda(t,s,r-q)|
\lesssim
\xi^\lambda(\xi-\eta)^{-\lambda},
\quad
|\p_y\Ecal_\lambda(t,s,r-q)|
\lesssim
\xi^{\lambda-1}(\xi-\eta)^{-\lambda}.
\end{equation}

\begin{equation}
\label{qq:46}
\begin{aligned}
&
\left|
\Ecal_\lambda(t,s,r-q)
-
\Ecal_\lambda(t,s,r+q)
\right|
\lesssim
\f{rqB_0^\lambda}{A_0(B_0-\eta)}
(\xi-\eta)^{-\lambda},\\
&
|\p_y\Ecal_\lambda(t,s,r-q)|
+
|\p_y\Ecal_\lambda(t,s,r+q)|
\lesssim
\f{(r+q)B_0^\lambda}{A_0(B_0-\eta)}
(\xi-\eta)^{-\lambda}.
\end{aligned}
\end{equation}

\begin{equation}
\label{eq:Ij}
I_j(\xi)
\lesssim
(1+\xi)^{p+j-1-\kappa-\lambda},
\qquad j=1,2.
\end{equation}

\begin{equation}\label{DELTA}
\xi^{j-\beta-\lambda}\log(2+\xi)
\lesssim
\xi^{j-\beta-\lambda}(1+\xi)^\delta
\lesssim
(1+\xi)^{p+j-1-\kappa-\lambda},
\end{equation}

\begin{proposition}
\label{prop:nonlinear}
Let $\kappa$ satisfy \eqref{eq:kappa-range}. Then, for any
$w\in X_\kappa$, $t\geq1$, and $0\leq r\leq t$,
\begin{equation}
\label{qq:47}
|\Tcal w(t,r)|
+
\f{1}{1+r}
\left|
\p_r(r\Tcal w)(t,r)
\right|
\lesssim
\|w\|_{X_\kappa}^{p}
\Acal(t,r)^{-1}\Bcal(t,r)^{1-\kappa}.
\end{equation}
In particular,
\begin{equation}
\label{qq:48}
\|\Tcal w\|_{X_\kappa}
\lesssim
\|w\|_{X_\kappa}^{p}.
\end{equation}
Moreover, for any $w,v\in X_\kappa$,
\begin{equation}
\label{qq:49}
\|\Tcal w-\Tcal v\|_{X_\kappa}
\lesssim
\big(
\|w\|_{X_\kappa}^{p-1}
+
\|v\|_{X_\kappa}^{p-1}
\big)
\|w-v\|_{X_\kappa}.
\end{equation}
\end{proposition}

\begin{proof}
Set
$
M=\|w\|_{X_\kappa}.
$
The support condition in $X_\kappa$ gives $G_w(s,q)=0$ for $q>s$.
For $0\leq q\leq s$, it follows from the definition of $X_\kappa$ in \eqref{eq:Xkappa}
that
\begin{equation}
\label{qq:50'}
|w(s,q)|
\leq
M(1+s+q)^{-1}(1+s-q)^{1-\kappa},
\end{equation}
and then
\begin{equation}
\label{qq:50}
|G_w(s,q)|
\lesssim
M^p q s^{-\beta}
(1+s+q)^{-p}
(1+s-q)^{-p(\kappa-1)}.
\end{equation}

We first estimate $r(\Tcal w)(t,r)$ for $r>0$. To this end, we consider the first
term on the right-hand side of \eqref{qq:40}
\[
\int_1^t
\int_{|t-s-r|}^{t-s+r}
\Ecal_\lambda(t,s,r-q)
G_w(s,q)\dd q\dd s.
\]
Under the change of variables in \eqref{qq:44}, the condition
$0\leq q\leq s$ implies
$
-\xi\leq\eta\leq0.
$
Moreover, recalling that $B_0=t-r$ and $A_0=t+r$ in \eqref{qq:44},
the condition $|t-s-r|\leq q\leq t-s+r$ gives
\[
\xi=s+q
\geq
s+|B_0-s|
\geq
B_0,
\quad
\xi\geq1,
\quad
\xi=s+q\leq A_0.
\]
Then
\[
\max\{1,B_0\}\leq\xi\leq A_0.
\]
Therefore, it follows from \eqref{qq:44}-\eqref{qq:45}, \eqref{qq:50} and \eqref{eq:Ij} that
\begin{equation}\label{QQ-44}
\begin{aligned}
\big|
\int_1^t
\int_{|t-s-r|}^{t-s+r}
\Ecal_\lambda(t,s,r-q)
G_w(s,q)\dd q\dd s
\big|
&\lesssim
M^p
\int_{\max\{1,B_0\}}^{A_0}
\xi^\lambda(1+\xi)^{-p}
I_1(\xi)\dd\xi\\
&\lesssim
M^p
\int_{B_0}^{A_0}
(1+\xi)^{-\kappa}\dd\xi\\
&=
\f{M^p}{\kappa-1}
\big[
\Bcal(t,r)^{1-\kappa}
-
\Acal(t,r)^{1-\kappa}
\big]\\
&=\f{M^p}{\kappa-1}\Bcal(t,r)^{1-\kappa}
\big[
1-
\big(
\f{\Bcal(t,r)}{\Acal(t,r)}
\big)^{\kappa-1}
\big]\\
&\lesssim
M^p r
\Acal(t,r)^{-1}
\Bcal(t,r)^{1-\kappa},
\end{aligned}
\end{equation}
where the last inequality follows from
$1-z^{\kappa-1}\lesssim1-z$ for $0\leq z\leq1$ and
$\Acal(t,r)-\Bcal(t,r)=2r$.

We next treat the remaining term on the right-hand side of
\eqref{qq:40}
\begin{equation}\label{YHC-020}
\begin{aligned}
\int_1^{t-r}
\int_0^{t-s-r}
\left[
\Ecal_\lambda(t,s,r-q)
-
\Ecal_\lambda(t,s,r+q)
\right]
G_w(s,q)\dd q\dd s.
\end{aligned}
\end{equation}
Note that this term vanishes when $B_0=t-r\leq1$. For $B_0>1$, the
integration region in \eqref{YHC-020} is
\[
1\leq s\leq B_0,
\quad
0\leq q\leq\min\{s,B_0-s\},
\]
which  imply
\[
1\leq\xi\leq B_0,
\quad
-\xi\leq\eta\leq0,\quad B_0-\eta\geq B_0.
\]
Then it follows from \eqref{qq:46},
\eqref{qq:50} and \eqref{eq:Ij} that
\begin{equation}\label{QQ-45}
\begin{aligned}
&\big|
\int_1^{t-r}
\int_0^{t-s-r}
\left[
\Ecal_\lambda(t,s,r-q)
-
\Ecal_\lambda(t,s,r+q)
\right]
G_w(s,q)\dd q\dd s
\big|\\
&\lesssim
M^p rA_0^{-1}B_0^{\lambda-1}
\int_1^{B_0}
(1+\xi)^{-p}I_2(\xi)\dd\xi\\
&\lesssim
M^p rA_0^{-1}B_0^{\lambda-1}
\int_1^{B_0}
(1+\xi)^{1-\kappa-\lambda}\dd\xi\\
&\lesssim
M^p rA_0^{-1}B_0^{1-\kappa}\\
&\lesssim
M^p r
\Acal(t,r)^{-1}
\Bcal(t,r)^{1-\kappa},
\end{aligned}
\end{equation}
where the third and last inequalities follow from
$\kappa+\lambda<2$ together with
$A_0\sim\Acal(t,r)$ and $B_0\sim\Bcal(t,r)$ for $B_0>1$. Combining \eqref{QQ-44} with
\eqref{QQ-45} yields
\begin{equation}
\label{qq:51}
|r(\Tcal w)(t,r)|
\lesssim
M^p r
\Acal(t,r)^{-1}
\Bcal(t,r)^{1-\kappa}.
\end{equation}

We next turn to deal with $\p_r(r\Tcal w)(t,r)$. From \eqref{qq:42}, we now estimate the first term of $\p_r(r\Tcal w)(t,r)$
\begin{equation}\label{YHC-08}
\begin{aligned}
4^\lambda
\int_1^t
\left[
G_w(s,t-s+r)
+
G_w(s,t-s-r)
\right]\dd s.
\end{aligned}
\end{equation}
For $G_w(s,t-s+r)$, the support condition
$t-s+r=A_0-s\leq s$ implies $s\geq\f{A_0}{2}$.
Setting $\theta=2s-A_0$, one has
\begin{equation}\label{QQ-46}
0\leq\theta\leq B_0,
\quad
1+s+(t-s+r)=\Acal(t,r),
\quad
1+s-(t-s+r)=1+\theta.
\end{equation}
Moreover, $q=t-s+r\leq A_0$ and $s\geq \f{A_0}{2}$ imply
\begin{equation}\label{QQ-46'}
q s^{-\beta}
\lesssim
A_0^{1-\beta}
\lesssim
\Acal(t,r)^{1-\beta}.
\end{equation}
Combining \eqref{QQ-45},  \eqref{QQ-46'} and \eqref{qq:50} yields
\begin{equation}\label{QQ-49}
\begin{aligned}
\int_1^t
|G_w(s,t-s+r)|\dd s
&\lesssim
M^p
\Acal(t,r)^{1-p-\beta}
\int_0^{B_0}
(1+\theta)^{-p(\kappa-1)}\dd\theta\lesssim
M^p\Acal(t,r)^{-\kappa}.
\end{aligned}
\end{equation}
Indeed, if $p(\kappa-1)<1$, then
\[
\int_0^{B_0}
(1+\theta)^{-p(\kappa-1)}\dd\theta
\lesssim
\Bcal(t,r)^{1-p(\kappa-1)}
\leq
\Acal(t,r)^{1-p(\kappa-1)}.
\]
Hence
\begin{equation}\label{YHC-09}
\begin{aligned}
M^p\Acal(t,r)^{1-p-\beta}
\int_0^{B_0}
(1+\theta)^{-p(\kappa-1)}\dd\theta
\lesssim M^p
\Acal(t,r)^{2-p-\beta-p(\kappa-1)}
\lesssim
M^p\Acal(t,r)^{-\kappa},
\end{aligned}
\end{equation}
where  we have used the fact of
$(p-1)\kappa>2-\beta$ in the last inequality.
If $p(\kappa-1)>1$, then the integral $\int_0^{B_0}
(1+\theta)^{-p(\kappa-1)}\dd\theta$ is uniformly bounded, and thus the condition
$\kappa<p+\beta-1$ implies \eqref{QQ-49}.
If $p(\kappa-1)=1$, as treated in \eqref{DELTA}, then \eqref{YHC-09} still holds.


We next treat the term which comes from \eqref{YHC-08}
\begin{equation}\label{YHC-011}
\begin{aligned}
\int_1^t
|G_w(s,t-s-r)|\dd s
=
\int_1^t
|G_w(s,B_0-s)|\dd s.
\end{aligned}
\end{equation}
First, it is assumed that $B_0>1$. Then
\begin{equation}\label{YHC-010}
\begin{aligned}
\int_1^t
|G_w(s,B_0-s)|\dd s
=
\int_1^{B_0}
|G_w(s,B_0-s)|\dd s
+
\int_{B_0}^{t}
|G_w(s,B_0-s)|\dd s .
\end{aligned}
\end{equation}
The first integral on the right-hand side of \eqref{YHC-010} follows from the estimate for
$G_w(s,t-s+r)$ by replacing $A_0$ with $B_0$. Thus,
\begin{equation}
\label{qq:54}
\int_1^{B_0}
|G_w(s,B_0-s)|\dd s
\lesssim
M^p\Bcal(t,r)^{-\kappa}.
\end{equation}
For the second integral on the right-hand side of \eqref{YHC-010}, by $s\geq B_0$ and the oddness of
$G_w(s,\cdot)$, one has
\[
|G_w(s,B_0-s)|
=|G_w(s,s-B_0)|.
\]
Using \eqref{qq:50} with
$q=s-B_0$ and $\xi=s+q=2s-B_0$, we arrive at
\begin{equation}
\label{qq:55}
\begin{aligned}
&\int_{B_0}^{t}
|G_w(s,B_0-s)|\dd s\lesssim
M^p(1+B_0)^{-p(\kappa-1)}
\int_{B_0}^{A_0}
(\xi+B_0)^{-\beta}
(1+\xi)^{-p}
(\xi-B_0)\dd\xi.
\end{aligned}
\end{equation}
In addition, by $B_0>1$ and $p+\beta>2$, it holds that
\begin{equation}
\label{qq:56}
\begin{aligned}
&\int_{B_0}^{A_0}
(\xi+B_0)^{-\beta}
(1+\xi)^{-p}
(\xi-B_0)\dd\xi\lesssim
\int_{B_0}^{2B_0}
B_0^{-\beta-p}
(\xi-B_0)\dd\xi
+\int_{2B_0}^{\infty}
\xi^{1-p-\beta}\dd\xi\lesssim
B_0^{2-p-\beta}.
\end{aligned}
\end{equation}
Then it follows from \eqref{qq:55} and the condition $(p-1)\kappa>2-\beta$ that
\begin{equation}
\label{qq:57}
\begin{aligned}
\int_{B_0}^{t}
|G_w(s,B_0-s)|\dd s
&\lesssim
M^p
\Bcal(t,r)^{-p(\kappa-1)}
B_0^{2-p-\beta}\lesssim
M^p\Bcal(t,r)^{-\kappa}.
\end{aligned}
\end{equation}
Combining \eqref{qq:54} and \eqref{qq:57} yields that for $B_0>1$,
\begin{equation}
\label{qq:58}
\int_1^t
|G_w(s,t-s-r)|\dd s
\lesssim
M^p\Bcal(t,r)^{-\kappa}.
\end{equation}
Next, consider the case of $B_0\leq1$ in \eqref{YHC-011}. Due to
$\Bcal(t,r)=1+B_0\sim1$,
we can arrive at $|B_0-s|=s-B_0$ for $s\geq1$. Substituting
$q=s-B_0$ into \eqref{qq:50} yields
\[
\begin{aligned}
|G_w(s,B_0-s)|
&\lesssim
M^p(s-B_0)s^{-\beta}
(1+2s-B_0)^{-p}
(1+B_0)^{-p(\kappa-1)}\lesssim
M^p(1+s)^{1-p-\beta}.
\end{aligned}
\]
Then the conditions $p+\beta>2$ and $\Bcal(t,r)\sim1$ imply
\[
\int_1^t
|G_w(s,B_0-s)|\dd s
\lesssim
M^p\int_1^\infty
(1+s)^{1-p-\beta}\dd s
\lesssim
M^p,
\]
which means
\begin{equation}\label{QQ-48}
\int_1^t
|G_w(s,t-s-r)|\dd s
\lesssim
M^p\Bcal(t,r)^{-\kappa}.
\end{equation}
Therefore, collecting the results in \eqref{QQ-49}, \eqref{qq:58} and  \eqref{QQ-48}, we conclude
\begin{equation}\label{QQ-50}
\begin{aligned}
4^\lambda \big|
\int_1^t
\left[
G_w(s,t-s+r)
+G_w(s,t-s-r)
\right]\dd s \big| 
&\lesssim
M^p
\left(
\Acal(t,r)^{-\kappa}
+\Bcal(t,r)^{-\kappa}
\right)\\
&\lesssim
M^p(1+r)
\Acal(t,r)^{-1}
\Bcal(t,r)^{1-\kappa},
\end{aligned}
\end{equation}
where in the last inequality we have used the facts of
$\Acal(t,r)=\Bcal(t,r)+2r$ and $\Bcal(t,r)\geq1$.



We next consider the second term on the right-hand side of
\eqref{qq:42}, namely,
\begin{equation}\label{YHC-013}
\begin{aligned}
\int_1^t
\int_{|t-s-r|}^{t-s+r}
\partial_y\Ecal_\lambda(t,s,r-q)
G_w(s,q)\dd q\dd s .
\end{aligned}
\end{equation}
Using the change of variables in \eqref{qq:44},  from the integration region in \eqref{YHC-013}, one has
\[
\xi=s+q\leq A_0,
\quad
\xi=s+q\geq s+|B_0-s|\geq B_0,
\]
while $\xi=s+q\geq1$. Hence,
\[
\max\{1,B_0\}\leq\xi\leq A_0.
\]
It follows from \eqref{qq:45} and \eqref{qq:50} that
\begin{equation}
\label{qq:62}
\begin{aligned}
&
\big|
\int_1^t
\int_{|t-s-r|}^{t-s+r}
\partial_y\Ecal_\lambda(t,s,r-q)
G_w(s,q)\dd q\dd s
\big|
\lesssim
M^p
\int_{\max\{1,B_0\}}^{A_0}
\xi^{\lambda-1}
(1+\xi)^{-p}
I_1(\xi)
\dd\xi .
\end{aligned}
\end{equation}
Applying \eqref{eq:Ij} with $j=1$ yields
\[
I_1(\xi)
\lesssim
(1+\xi)^{p-\kappa-\lambda}.
\]
Therefore,
\begin{equation}
\label{qq:63}
\begin{aligned}
&
\big|
\int_1^t
\int_{|t-s-r|}^{t-s+r}
\partial_y\Ecal_\lambda(t,s,r-q)
G_w(s,q)\dd q\dd s
\big|
\lesssim
M^p
\int_{B_0}^{A_0}
(1+\xi)^{-\kappa-1}
\dd\xi .
\end{aligned}
\end{equation}
Similarly treated as in \eqref{QQ-44}, one has
\begin{equation}
\label{qq:64}
\begin{aligned}
&
\big|
\int_1^t
\int_{|t-s-r|}^{t-s+r}
\partial_y\Ecal_\lambda(t,s,r-q)
G_w(s,q)\dd q\dd s
\big|
\lesssim
M^p
r\Acal(t,r)^{-1}
\Bcal(t,r)^{1-\kappa}.
\end{aligned}
\end{equation}

It remains to consider the last term on the right-hand side of
\eqref{qq:42}
\[
\int_1^{t-r}
\int_0^{t-s-r}
\left[
\partial_y\Ecal_\lambda(t,s,r-q)
-
\partial_y\Ecal_\lambda(t,s,r+q)
\right]
G_w(s,q)\dd q\dd s,
\]
which vanishes when $B_0\leq1$. For $B_0>1$, by the same
change of variables as in the proof of \eqref{QQ-45}, together
with \eqref{qq:46}, \eqref{qq:50} and \eqref{eq:Ij}, we have
\begin{equation}
\label{qq:64'}
\begin{aligned}
&
\big|
\int_1^{t-r}
\int_0^{t-s-r}
\left[
\partial_y\Ecal_\lambda(t,s,r-q)
-
\partial_y\Ecal_\lambda(t,s,r+q)
\right]
G_w(s,q)\dd q\dd s
\big|\\
&\lesssim
M^p
A_0^{-1}B_0^{\lambda-1}
\big[
r
\int_1^{B_0}
(1+\xi)^{-p}
I_1(\xi)\dd\xi
+
\int_1^{B_0}
(1+\xi)^{-p}
I_2(\xi)\dd\xi
\big].
\end{aligned}
\end{equation}
By \eqref{eq:Ij}, one can obtain
\[
\begin{aligned}
&
r
\int_1^{B_0}
(1+\xi)^{-p}
I_1(\xi)\dd\xi
+
\int_1^{B_0}
(1+\xi)^{-p}
I_2(\xi)\dd\xi\lesssim
r
\int_1^{B_0}
(1+\xi)^{-\kappa-\lambda}
\dd\xi
+
\int_1^{B_0}
(1+\xi)^{1-\kappa-\lambda}
\dd\xi .
\end{aligned}
\]
Since
$
1<\kappa+\lambda<2,
$
it follows that
\[
\begin{aligned}
&
r
\int_1^{B_0}
(1+\xi)^{-\kappa-\lambda}
\dd\xi
+
\int_1^{B_0}
(1+\xi)^{1-\kappa-\lambda}
\dd\xi\lesssim
r+B_0^{2-\kappa-\lambda}.
\end{aligned}
\]
Then
\begin{equation}
\label{qq:65}
\begin{aligned}
&
\big|
\int_1^{t-r}
\int_0^{t-s-r}
\left[
\partial_y\Ecal_\lambda(t,s,r-q)
-
\partial_y\Ecal_\lambda(t,s,r+q)
\right]
G_w(s,q)\dd q\dd s
\big|\\
&\lesssim
M^pA_0^{-1}B_0^{\lambda-1}
\big(
r+B_0^{2-\kappa-\lambda}
\big)\\
&\lesssim
M^pA_0^{-1}
\big(
rB_0^{\lambda-1}
+B_0^{1-\kappa}
\big).
\end{aligned}
\end{equation}
Due to $B_0>1$, one has
$A_0\sim\Acal(t,r)$ and $B_0\sim\Bcal(t,r)$.
Then it follows from the condition $\kappa<2-\lambda$  that
\begin{equation}
\label{qq:66}
\begin{aligned}
&
\big|
\int_1^{t-r}
\int_0^{t-s-r}
\left[
\partial_y\Ecal_\lambda(t,s,r-q)
-\partial_y\Ecal_\lambda(t,s,r+q)
\right]
G_w(s,q)\dd q\dd s
\big|\\
&\lesssim
M^p\Acal(t,r)^{-1}
\left(
r\Bcal(t,r)^{1-\kappa}
+\Bcal(t,r)^{1-\kappa}
\right)\\
&\lesssim
M^p(1+r)
\Acal(t,r)^{-1}
\Bcal(t,r)^{1-\kappa}.
\end{aligned}
\end{equation}
By \eqref{QQ-50}, \eqref{qq:64} and
\eqref{qq:66}, we can conclude
\begin{equation}
\label{qq:67}
\left|
\partial_r(r\Tcal w)(t,r)
\right|
\lesssim
M^p(1+r)
\Acal(t,r)^{-1}
\Bcal(t,r)^{1-\kappa}.
\end{equation}
Combining \eqref{qq:51} and \eqref{qq:67} yields that for $r>0$,
\begin{equation}
\label{qq:68}
|\Tcal w(t,r)|
+\frac1{1+r}
\left|
\partial_r(r\Tcal w)(t,r)
\right|
\lesssim
M^p
\Acal(t,r)^{-1}
\Bcal(t,r)^{1-\kappa}.
\end{equation}
For $r=0$, the oddness of $r\Tcal w(t,r)$ with respect to $r$
implies
\[
\Tcal w(t,0)
=\lim_{r\to0}
\frac{r\Tcal w(t,r)}{r}
=\partial_r(r\Tcal w)(t,0).
\]
The estimate at $r=0$ follows from \eqref{qq:67} by letting
$r$ approach zero.
Hence, \eqref{qq:47} holds for $0\leq r\leq t$.
Then \eqref{qq:48} follows immediately from the definition of
$X_\kappa$ and \eqref{qq:47}.

It remains to prove \eqref{qq:49}. Notice that for any $w,v\in X_\kappa$, one has
\[
\bigl||w|^p-|v|^p\bigr|
\lesssim
|w-v|
\left(
|w|^{p-1}
+|v|^{p-1}
\right).
\]
Applying \eqref{qq:50'} to $w$, $v$, and $w-v$, respectively,
one obtains
\[
\begin{aligned}
|G_w(s,q)-G_v(s,q)|
&\lesssim
q s^{-\beta}
\bigl||w(s,q)|^p-|v(s,q)|^p\bigr|
\\
&\lesssim
\big(\|w\|_{X_\kappa}^{p-1}
+\|v\|_{X_\kappa}^{p-1}
\big)
\|w-v\|_{X_\kappa}
q s^{-\beta}
(1+s+q)^{-p}
(1+s-q)^{-p(\kappa-1)}.
\end{aligned}
\]
Since the above expression has the same form as \eqref{qq:50},
the estimates derived above remain valid with $G_w$ replaced by
$G_w-G_v$.
Therefore, \eqref{qq:49}  follows immediately from the resulting
estimate and the definition of $X_\kappa$. This completes the proof.
\end{proof}

\end{document}
